% Unit 5 exam -- 25 MCQ + 1 short + 1 long FRQ = 39 pts, 69 min.
\documentclass{apchem-exam}

\apcunit{5}
\apcunittitle{Kinetics}
\apcdoctitle{Unit Exam}
\apcsubtitle{CED 5.1--5.11 \textbullet\ Fri Nov 20 \textbullet\ 69 minutes \textbullet\ 39 points}

\begin{document}
\apctitleblock

\begin{apcnote}[Format]
Section I: 25 multiple choice, \textbf{no calculator}, 37 minutes.
Section II: one short and one long free response, calculator permitted,
32 minutes. Constants and equations are provided with Section II.
\end{apcnote}

\examsection{I}{Multiple Choice}{25 questions \textbullet\ 37 minutes \textbullet\ 1 point each}{\nocalculator}

% ---- rates -----------------------------------------------------------------
\begin{mcq}
  For $\ce{2A -> B}$, if A disappears at \SI{0.040}{\Molar\per\second}, the
  rate of the reaction is
  \choices
    \choice{\SI{0.080}{\Molar\per\second}}
    \correct{\SI{0.020}{\Molar\per\second}}
    \choice{\SI{0.040}{\Molar\per\second}}
    \choice{\SI{0.010}{\Molar\per\second}}
\end{mcq}
\why{Divide by the coefficient of 2.}

\begin{mcq}
  For $\ce{N2 + 3H2 -> 2NH3}$, \ce{H2} is consumed
  \choices
    \correct{three times as fast as \ce{N2}}
    \choice{at the same rate as \ce{N2}}
    \choice{twice as fast as \ce{N2}}
    \choice{half as fast as \ce{N2}}
\end{mcq}
\why{The coefficient ratio is 3:1.}

\begin{mcq}
  The instantaneous rate at a given moment is
  \choices
    \correct{the slope of the tangent to the concentration--time curve}
    \choice{the slope of a line joining two points on the curve}
    \choice{the average of the initial and final rates}
    \choice{always equal to the initial rate}
\end{mcq}
\why{The average rate is the secant slope; the instantaneous rate is the
tangent.}

\begin{mcq}
  Rate laws are determined from initial rates because at $t = 0$
  \choices
    \correct{no products have accumulated to complicate the measurement}
    \choice{the temperature is lowest}
    \choice{the rate constant is largest}
    \choice{the reaction is at equilibrium}
\end{mcq}
\why{No reverse reaction and no product interference.}

% ---- rate laws -------------------------------------------------------------
\begin{mcq}
  In the rate law rate $= k[\ce{A}]^m[\ce{B}]^n$, the exponents $m$ and $n$
  \choices
    \choice{equal the coefficients in the balanced equation}
    \correct{must be determined experimentally}
    \choice{are always whole numbers equal to 1}
    \choice{depend on the concentrations used}
\end{mcq}
\why{Coefficients give orders only for an elementary step.}
\distractor{A}{$\ce{2N2O5 -> 4NO2 + O2}$ is first order --- the standard
counterexample}

\begin{mcq}
  Doubling $[\ce{A}]$ leaves the rate unchanged. The order in A is
  \choices
    \correct{0}
    \choice{1}
    \choice{2}
    \choice{$\tfrac{1}{2}$}
\end{mcq}
\why{$2^0 = 1$.}

\begin{mcq}
  Tripling $[\ce{A}]$ multiplies the rate by 9. The order in A is
  \choices
    \choice{1}
    \correct{2}
    \choice{3}
    \choice{9}
\end{mcq}
\why{$3^2 = 9$.}

\begin{mcq}
  A rate constant has units of \si{\per\Molar\per\second}. The overall
  order is
  \choices
    \choice{0}
    \choice{1}
    \correct{2}
    \choice{3}
\end{mcq}
\why{Units of $k$ are fixed by the overall order.}

\begin{mcq}
  For rate $= k[\ce{A}]^2[\ce{B}]$, tripling $[\ce{A}]$ and halving
  $[\ce{B}]$ changes the rate by a factor of
  \choices
    \choice{1.5}
    \choice{3}
    \correct{4.5}
    \choice{9}
\end{mcq}
\why{$3^2 \times \tfrac{1}{2}$.}

\begin{mcq}
  The rate constant $k$ is independent of
  \choices
    \correct{concentration}
    \choice{temperature}
    \choice{activation energy}
    \choice{the presence of a catalyst}
\end{mcq}
\why{$k$ changes with temperature, $E_a$ and catalysis, but not with how
much you mix.}

% ---- integrated -------------------------------------------------------------
\begin{mcq}
  A plot of $\ln[\ce{A}]$ against time is linear. The reaction is
  \choices
    \choice{zero order}
    \correct{first order}
    \choice{second order}
    \choice{third order}
\end{mcq}
\why{Each order linearizes a different function.}

\begin{mcq}
  A plot of $1/[\ce{A}]$ against time is linear with a positive slope. The
  slope equals
  \choices
    \correct{$k$}
    \choice{$-k$}
    \choice{$1/k$}
    \choice{$0.693/k$}
\end{mcq}
\why{$1/[\ce{A}] = kt + 1/[\ce{A}]_0$.}

\begin{mcq}
  For which order is the half-life independent of the starting
  concentration?
  \choices
    \choice{zero}
    \correct{first}
    \choice{second}
    \choice{all of them}
\end{mcq}
\why{$t_{1/2} = 0.693/k$ contains no $[\ce{A}]_0$.}

\begin{mcq}
  A first-order reaction has $k = \SI{0.0693}{\per\second}$. Its half-life
  is
  \choices
    \choice{\SI{0.693}{\second}}
    \choice{\SI{6.93}{\second}}
    \correct{\SI{10.0}{\second}}
    \choice{\SI{100}{\second}}
\end{mcq}
\why{$0.693/0.0693$.}

\begin{mcq}
  An isotope has a half-life of \SI{8.0}{\day}. Starting from
  \SI{40.0}{\milli\gram}, the mass after \SI{24}{\day} is
  \choices
    \choice{\SI{20.0}{\milli\gram}}
    \choice{\SI{10.0}{\milli\gram}}
    \correct{\SI{5.0}{\milli\gram}}
    \choice{\SI{2.5}{\milli\gram}}
\end{mcq}
\why{Three half-lives: $40 \to 20 \to 10 \to 5$.}

\begin{mcq}
  For a second-order reaction, as the reaction proceeds the successive
  half-lives
  \choices
    \correct{get longer}
    \choice{get shorter}
    \choice{stay the same}
    \choice{become zero}
\end{mcq}
\why{$t_{1/2} = 1/(k[\ce{A}]_0)$ and the concentration is falling.}

% ---- mechanisms -------------------------------------------------------------
\begin{mcq}
  For the elementary step $\ce{2NO -> N2O2}$, the rate law is
  \choices
    \choice{rate $= k[\ce{NO}]$}
    \correct{rate $= k[\ce{NO}]^2$}
    \choice{rate $= k[\ce{N2O2}]$}
    \choice{cannot be determined without data}
\end{mcq}
\why{For an elementary step --- and only then --- molecularity gives the
order.}

\begin{mcq}
  A species produced in one step of a mechanism and consumed in a later one
  is
  \choices
    \correct{an intermediate}
    \choice{a catalyst}
    \choice{a transition state}
    \choice{a spectator}
\end{mcq}
\why{A catalyst is consumed first and regenerated later --- the reverse
order.}

\begin{mcq}
  A proposed mechanism is acceptable only if
  \choices
    \correct{its steps sum to the overall equation and it predicts the
      observed rate law}
    \choice{it contains no intermediates}
    \choice{every step is bimolecular}
    \choice{the first step is always slow}
\end{mcq}
\why{Both conditions, and even then it is consistent rather than proven.}

\begin{mcq}
  Termolecular elementary steps are rare because
  \choices
    \correct{a simultaneous three-particle collision is highly improbable}
    \choice{they violate conservation of mass}
    \choice{they always have zero activation energy}
    \choice{they cannot be rate-determining}
\end{mcq}
\why{Probability, not principle.}

% ---- collision, profiles, catalysis -----------------------------------------
\begin{mcq}
  A collision leads to reaction only if the particles collide with
  sufficient energy AND
  \choices
    \correct{the correct orientation}
    \choice{equal masses}
    \choice{opposite charges}
    \choice{the same speed}
\end{mcq}
\why{Energy and orientation are both required.}

\begin{mcq}
  Raising the temperature increases the rate mainly because
  \choices
    \correct{a larger fraction of collisions exceeds the activation energy}
    \choice{collisions become much more frequent}
    \choice{the activation energy decreases}
    \choice{$\Delta H$ becomes more negative}
\end{mcq}
\why{The frequency effect is small; the distribution effect dominates.}
\distractor{C}{only a catalyst changes $E_a$}

\begin{mcq}
  On an energy profile, the activation energy is measured from
  \choices
    \correct{the reactants up to the transition state}
    \choice{the reactants to the products}
    \choice{the products up to the transition state}
    \choice{the baseline to the products}
\end{mcq}
\why{That last one from products would be the reverse $E_a$.}

\begin{mcq}
  A reaction has $E_a(\text{forward}) = \SI{80}{\kilo\joule}$ and
  $E_a(\text{reverse}) = \SI{120}{\kilo\joule}$. Then $\Delta H$ is
  \choices
    \correct{\SI{-40}{\kilo\joule}}
    \choice{\SI{+40}{\kilo\joule}}
    \choice{\SI{+200}{\kilo\joule}}
    \choice{\SI{-200}{\kilo\joule}}
\end{mcq}
\why{$E_a(\text{fwd}) - E_a(\text{rev})$.}

\begin{mcq}
  Adding a catalyst to a reaction at equilibrium
  \choices
    \correct{speeds up both directions equally and does not shift the
      equilibrium}
    \choice{shifts the equilibrium toward products}
    \choice{increases $K$}
    \choice{makes $\Delta H$ more negative}
\end{mcq}
\why{It lowers $E_a$ for forward and reverse by the same amount.}

\examsection{II}{Free Response}{2 questions \textbullet\ 32 minutes}{\calculatorok}

\begin{apcnote}[Data]
rate $= k[\ce{A}]^m[\ce{B}]^n$ \textbullet\
$\ln[\ce{A}] = -kt + \ln[\ce{A}]_0$ \textbullet\
$1/[\ce{A}] = kt + 1/[\ce{A}]_0$ \textbullet\
$t_{1/2} = 0.693/k$ (first order) \textbullet\
$k = Ae^{-E_a/RT}$ \textbullet\
$\ln(k_2/k_1) = -\tfrac{E_a}{R}(\tfrac{1}{T_2}-\tfrac{1}{T_1})$
\textbullet\ $R = \SI{8.314}{\joule\per\mole\per\kelvin}$.
\end{apcnote}

% ---- FRQ 1 (short) ---------------------------------------------------------
\frq{short}{4}{CED 5.3 \textbullet\ integrated rate laws}

A reaction is followed and three plots are prepared. Only the plot of
$1/[\ce{A}]$ against time is linear; it has slope $+0.42$ and intercept
$5.0$.

\begin{parts}
  \item State the order of the reaction and justify from the graphs.
        \answerlines[2]{Second order. Each order linearizes a different
        function, and here it is $1/[\ce{A}]$ against $t$ that is straight
        --- which is the second-order integrated form.}
  \item Give $k$ with units and $[\ce{A}]_0$.
        \answerlines[2]{Slope $= k = \SI{0.42}{\per\Molar\per\second}$;
        intercept $= 1/[\ce{A}]_0 = 5.0$, so
        $[\ce{A}]_0 = \SI{0.20}{\Molar}$.}
  \item Calculate the half-life at the start of the reaction.
        \workspace[2cm]
        \keyonly{\begin{apcanswer}$t_{1/2} = 1/(k[\ce{A}]_0) =
        1/[(0.42)(0.20)] = \mathbf{12}~\si{\second}$\end{apcanswer}}
\end{parts}

\begin{rubric}
  \rpoint{1}{(a) second order, justified by naming which plot is linear ---
    ``the graph is linear'' alone earns 0}
  \rpoint{2}{(b) 1 pt $k = 0.42$ with units; 1 pt
    $[\ce{A}]_0 = \SI{0.20}{\Molar}$ from the intercept}
  \rpoint{1}{(c) $t_{1/2} = 12$~s using the second-order expression ---
    using $0.693/k$ earns 0}
\end{rubric}

% ---- FRQ 2 (long) ----------------------------------------------------------
\frq{long}{10}{CED 5.2, 5.5, 5.8, 5.9 \textbullet\ a complete kinetics analysis}

The reaction $\ce{2NO(g) + O2(g) -> 2NO2(g)}$ was studied at
\SI{300}{\kelvin}.

\begin{center}
\begin{tabular}{@{}cccc@{}}
\hline
\textbf{Exp} & $[\ce{NO}]_0$ (M) & $[\ce{O2}]_0$ (M) &
  \textbf{initial rate} (\si{\Molar\per\second}) \\
\hline
1 & 0.010 & 0.010 & $2.5\times10^{-5}$ \\
2 & 0.020 & 0.010 & $1.0\times10^{-4}$ \\
3 & 0.010 & 0.020 & $5.0\times10^{-5}$ \\
\hline
\end{tabular}
\end{center}

\begin{parts}
  \item Determine the order with respect to each reactant, citing the
        experiments compared. \answerlines[3]{Comparing 1 and 2, only
        $[\ce{NO}]$ changes: doubling it multiplies the rate by 4, so the
        order in \ce{NO} is 2. Comparing 1 and 3, doubling $[\ce{O2}]$
        doubles the rate, so the order in \ce{O2} is 1.}
  \item Write the rate law and calculate $k$ with units.
        \workspace[2.4cm]
        \keyonly{\begin{apcanswer}rate $= k[\ce{NO}]^2[\ce{O2}]$;
        $k = 2.5\times10^{-5}/[(0.010)^2(0.010)] =
        \mathbf{25}~\si{\per\Molar\squared\per\second}$\end{apcanswer}}
  \item Predict the initial rate when $[\ce{NO}] = 0.030$~M and
        $[\ce{O2}] = 0.015$~M. \workspace[2.2cm]
        \keyonly{\begin{apcanswer}$(25)(0.030)^2(0.015) =
        \mathbf{3.4\times10^{-4}}~\si{\Molar\per\second}$\end{apcanswer}}
  \item The following mechanism is proposed:
        step 1 (fast equilibrium) $\ce{2NO <=> N2O2}$;
        step 2 (slow) $\ce{N2O2 + O2 -> 2NO2}$.
        Show that it predicts the observed rate law.
        \workspace[3cm]
        \keyonly{\begin{apcanswer}The slow step gives
        rate $= k_2[\ce{N2O2}][\ce{O2}]$, but \ce{N2O2} is an intermediate.
        From step 1's equilibrium, $K_1 = [\ce{N2O2}]/[\ce{NO}]^2$, so
        $[\ce{N2O2}] = K_1[\ce{NO}]^2$. Substituting gives
        rate $= k_2K_1[\ce{NO}]^2[\ce{O2}]$, which matches the observed form
        with $k = k_2K_1$.\end{apcanswer}}
  \item Identify the intermediate, and explain why it must not appear in
        the final rate law. \answerlines[2]{\ce{N2O2}. A rate law must be
        written in terms of species whose concentrations can be measured
        and controlled; an intermediate is never mixed in and its
        concentration is not independently known.}
  \item When the temperature is raised to \SI{310}{\kelvin}, $k$ increases
        by a factor of 1.9. Calculate $E_a$. \workspace[2.8cm]
        \keyonly{\begin{apcanswer}$\ln(1.9) = 0.642 =
        -(E_a/8.314)(1/310 - 1/300)$. The bracket is
        $-1.075\times10^{-4}$, so
        $E_a = (0.642)(8.314)/(1.075\times10^{-4}) =
        \mathbf{5.0\times10^{4}}~\si{\joule\per\mole} =
        \SI{50}{\kilo\joule\per\mole}$\end{apcanswer}}
\end{parts}

\begin{rubric}
  \rpoint{2}{(a) 1 pt both orders correct; 1 pt cites the specific
    experiments compared for each --- orders with no justification earn 1}
  \rpoint{2}{(b) 1 pt rate law; 1 pt $k = 25$ \emph{with units}}
  \rpoint{1}{(c) $3.4\times10^{-4}$~\si{\Molar\per\second}}
  \rpoint{3}{(d) 1 pt writes the slow step's rate law; 1 pt uses $K_1$ to
    express $[\ce{N2O2}]$; 1 pt substitutes to reach the observed form ---
    an answer still containing \ce{N2O2} earns at most 1}
  \rpoint{1}{(e) \ce{N2O2} identified with a reason referring to
    measurability}
  \rpoint{1}{(f) $E_a \approx \SI{50}{\kilo\joule\per\mole}$}
\end{rubric}

\examtotal

\end{document}
